%% ============================================================
%%  SECCIÓN III — Objetivos
%%  Responsable: Minilux
%% ============================================================

\section{Objetivos}

\subsection{Objetivo general}

Determinar experimentalmente el coeficiente de tensión superficial del agua
mediante dos métodos independientes basados en el equilibrio estático (torques
y fuerzas), comparando los resultados obtenidos con el valor de referencia
bibliográfico.

\subsection{Objetivos específicos}

\begin{itemize}
    \item Medir las dimensiones geométricas del anillo metálico (radios interno
    y externo) y determinar las masas de los jinetillos empleados en la balanza de
    Mohr--Westphal.

    \item Determinar la fuerza de tensión superficial a partir del equilibrio
    de torques en la balanza, utilizando cinco configuraciones distintas de
    jinetillos sobre el brazo mayor.

    \item Calcular el coeficiente de tensión superficial $\gamma$ para cada
    medición del primer método y obtener su valor promedio junto con su
    incertidumbre asociada mediante la desviación estándar de la media.

    \item Medir las dimensiones geométricas de la película jabonosa ($a$, $b$,
    $h$) y determinar la masa y el peso del tubo inferior en el segundo método.

    \item Demostrar analíticamente la fórmula de equilibrio mecánico para el
    segundo método y calcular el coeficiente $\gamma$ a partir de ella.

    \item Comparar los valores de $\gamma$ obtenidos por ambos métodos con el
    valor aceptado para el agua pura a \SI{20}{\celsius}
    ($\gamma_{\text{ref}} = \SI{72.8e-3}{\newton\per\meter}$), cuantificando
    el error relativo porcentual en cada caso.

    \item Identificar las fuentes de error experimental y analizar las discrepancias
    entre ambos métodos, en particular el efecto del surfactante sobre las fuerzas
    cohesivas del fluido.
\end{itemize}